\documentclass[UTF8,11pt]{ctexart}
\usepackage[a4paper,margin=22mm]{geometry}
\usepackage{graphicx,enumitem,amssymb,tabularx}
\usepackage[normalem]{ulem}
\usepackage{xeCJKfntef}
\setlength{\parindent}{0pt}
\setlength{\parskip}{.6em}
\sloppy
\begin{document}
\section*{2025年6月英语六级真题(第1套)}
\subsection*{2025 年6 月大学英语六级考试真题（第1 套）}


\subsection*{\textbf{Part I Writing} \textbf{(30 minutes)}}


\textbf{Directions:}\textit{ For this part, you are allowed 30 minutes to write an essay that begins with the sentence }\textbf{\textit{“As social media}} \textbf{\textit{is used more and more extensively, there is a}} \textbf{\textit{growing awareness of the importance of using}} \textbf{\textit{it properly and}} \textbf{\textit{responsibly.”}}\textit{ You can make comments, cite examples or use your personal experiences to develop your essay. You} \textit{should write at least }\textit{\uline{150 }}\textit{words but no more than }\textit{\uline{200 }}\textit{words.}


\textbf{\textit{\uline{You should copy the sentence given in quotes at the beginning of your essay.}}}


\subsection*{\textbf{Part II} \textbf{Listening Comprehension} \textbf{(30 minutes)}}


\subsection*{\textbf{Section A}}


\textbf{Directions:} \textit{In this section, you will hear two long conversations. At the end of each conversation, you will hear four} \textit{questions. Both the conversation and the questions will be spoken only once. After you hear a question, you must choose} \textit{the best answer from the four choices marked A, B, C and D. Then mark the corresponding letter on }\textbf{\textit{Answer Sheet 1}} \textit{with a single line through the centre.}


\textbf{Questions 1 to 4 are based on the conversation you have just heard.}


1.


\begin{itemize}[leftmargin=2em,itemsep=.2em,topsep=.4em]\item[A.] Met the computer technician.
\item[B.] Told the man about her trouble.
\item[C.] Called the man’s company.
\item[D.] Visited Alpha Maintenance.
\end{itemize}

2.


\begin{itemize}[leftmargin=2em,itemsep=.2em,topsep=.4em]\item[A.] Consulted someone in charge.
\item[B.] Came as soon as possible.
\item[C.] Informed the central office at once.
\item[D.] Designated an engineer to the case.
\end{itemize}

\noindent\begin{tabularx}{\linewidth}{@{}lXXXX@{}}3. & \textbf{A.} Frustration. & \textbf{B.} Intimidation. & \textbf{C.} Desperation. & \textbf{D.} Indignation.\\\end{tabularx}\par

4.


\begin{itemize}[leftmargin=2em,itemsep=.2em,topsep=.4em]\item[A.] Whether the contents have been backed up.
\item[B.] Whether they can find help somewhere else.
\item[C.] Whether all data stored on the hard drive has got lost.
\item[D.] Whether they need to wipe the system directories clean.
\end{itemize}

\textbf{Questions 5 to 8 are based on the conversation you have just heard.}


\noindent\begin{tabularx}{\linewidth}{@{}lXXXX@{}}5. & \textbf{A.} It’s boring. & \textbf{B.} It’s challenging. & \textbf{C.} It’s a beautiful thing. & \textbf{D.} It’s unlike most jobs.\\\end{tabularx}\par

6.


\begin{itemize}[leftmargin=2em,itemsep=.2em,topsep=.4em]\item[A.] Arbitrating between disagreeing solicitors.
\item[B.] Preventing disputes from escalating.
\item[C.] Buying and selling property.
\item[D.] Mediating land disputes.
\end{itemize}

7.


\begin{itemize}[leftmargin=2em,itemsep=.2em,topsep=.4em]\item[A.] Courts are intended for fixing major legal disputes.
\item[B.] Courts are getting too bureaucratic to function.
\item[C.] Courts can be frustrating and expensive.
\item[D.] Courts can be frightening and arbitrary.
\end{itemize}

8.


\begin{itemize}[leftmargin=2em,itemsep=.2em,topsep=.4em]\item[A.] The ability to make arguments in a unique way.
\item[B.] The skill of preventing conflicts between parties concerned.
\item[C.] The skill of foreseeing any potential stakes in their work.
\item[D.] The ability to express themselves clearly and forcefully.
\end{itemize}

\subsection*{\textbf{Section B}}


\textbf{Directions:} \textit{In this section, you will hear two passages. At the end of each passage, you will hear three or four questions.} \textit{Both the passage and the questions will be spoken only once. After you hear a question, you must choose the best answer} \textit{from the four choices marked A, B, C and D. Then mark the corresponding letter on }\textbf{\textit{Answer Sheet 1}}\textit{ with a single line} \textit{through the centre.}


\textbf{Questions 9 to 11 are based on the passage you have just heard.}


9.


\begin{itemize}[leftmargin=2em,itemsep=.2em,topsep=.4em]\item[A.] They looked into the relationship between one’s prior knowledge and creativity.
\end{itemize}

\begin{itemize}[leftmargin=2em,itemsep=.2em,topsep=.4em]\item[B.] They offered novel ways to help older adults to keep their memories from aging.
\item[C.] They proposed an explanation for old people’s difficulty in retrieving memories.
\item[D.] They advanced a new model concerning human information storage and retrieval.
\end{itemize}

10.


\begin{itemize}[leftmargin=2em,itemsep=.2em,topsep=.4em]\item[A.] Young adults rely on memory while older adults keep notes as a reminder.
\item[B.] Older adults often retrieve irrelevant memories along with what they want.
\item[C.] Young adults accumulate knowledge much more quickly than older adults.
\item[D.] Older adults generally perform cognitive tasks much slower than young adults.
\end{itemize}

11.


\begin{itemize}[leftmargin=2em,itemsep=.2em,topsep=.4em]\item[A.] They show preserved, and sometimes enhanced, creativity.
\item[B.] They frequently suffer from disorderly crowded memories.
\item[C.] They can rely on their accumulated wisdom in an emergency.
\item[D.] They may well be served by forgetting their prior knowledge.
\end{itemize}

\textbf{Questions 12 to 15 are based on the passage you have just heard.}


12.


\begin{itemize}[leftmargin=2em,itemsep=.2em,topsep=.4em]\item[A.] They are actually proud of the goods and services they produce.
\item[B.] They are worried about being alienated from the outside world.
\item[C.] They are rarely in control of many things related to their work.
\item[D.] They are longing to share the profits made from their work.
\end{itemize}

13.


\begin{itemize}[leftmargin=2em,itemsep=.2em,topsep=.4em]\item[A.] The steady decrease in productivity.
\item[B.] The adverse effect on physical health.
\item[C.] The feeling of being time-poor.
\item[D.] The rising economic inequality.
\end{itemize}

14.


\begin{itemize}[leftmargin=2em,itemsep=.2em,topsep=.4em]\item[A.] It alters the structure of work.
\item[B.] It puts jobs and wages at risk.
\item[C.] It liberates people from tedious and laborious work.
\item[D.] It creates new work opportunities in the IT industry.
\end{itemize}

15.


\begin{itemize}[leftmargin=2em,itemsep=.2em,topsep=.4em]\item[A.] Finding meaning in work.
\item[B.] Prioritizing life over work.
\item[C.] Improving relationships in the community.
\item[D.] Realizing one’s social value in the workplace.
\end{itemize}

\subsection*{\textbf{Section C}}


\textbf{Directions:}\textit{ In this section, you will hear three recordings of lectures or talks followed by three or four questions. The} \textit{recordings will be played only once. After you hear a question, you must choose the best answer from the four choices} \textit{marked A, B, C and D. Then mark the corresponding letter on }\textbf{\textit{Answer Sheet 1}}\textit{ with a single line through the centre.}


\textbf{Questions 16 to 18 are based on the recording you have just heard.}


16.


\begin{itemize}[leftmargin=2em,itemsep=.2em,topsep=.4em]\item[A.] Affect our attitude to novel tasks.
\item[B.] Distinguish us in the social world.
\item[C.] Outweigh IQ in importance.
\item[D.] Impact academic achievements.
\end{itemize}

17.


\begin{itemize}[leftmargin=2em,itemsep=.2em,topsep=.4em]\item[A.] By pressing a hidden button.
\item[B.] By pushing a big button on top.
\item[C.] By pressing two buttons at the same time.
\item[D.] By helping the babies push the right button.
\end{itemize}

18.


\begin{itemize}[leftmargin=2em,itemsep=.2em,topsep=.4em]\item[A.] Perform difficult tasks successfully just by observing how adults acted.
\item[B.] Make generalizable inferences about persistence from a few examples.
\item[C.] Adapt themselves to different social contexts.
\item[D.] Work hard to interact with experimenters.
\end{itemize}

\textbf{Questions 19 to 21 are based on the recording you have just heard.}


19.


\begin{itemize}[leftmargin=2em,itemsep=.2em,topsep=.4em]\item[A.] Offering advice on overcoming habitual lateness.
\item[B.] Exemplifying various recreational opportunities.
\item[C.] Scrutinizing individuals’ defining traits.
\item[D.] Suggesting ways for setting priorities.
\end{itemize}

20.


\begin{itemize}[leftmargin=2em,itemsep=.2em,topsep=.4em]\item[A.] Make the breakfast simpler.
\item[B.] Take the alarms seriously.
\item[C.] Ready yourself in the early dawn.
\item[D.] Get prepared the night before.
\end{itemize}

21.


\begin{itemize}[leftmargin=2em,itemsep=.2em,topsep=.4em]\item[A.] Finish the prior task 30 minutes earlier.
\item[B.] Keep ourselves from hitting a bump.
\item[C.] Leave time in between activities.
\item[D.] Try to avoid possible hold-ups.
\end{itemize}

\textbf{Questions 22 to 25 are based on the recording you have just heard.}


\noindent\begin{tabularx}{\linewidth}{@{}lXXXX@{}}22. & \textbf{A.} Their lifestyles vary. & \textbf{B.} Their traits vary. & \textbf{C.} They have different customs. & \textbf{D.} They have different feels.\\\end{tabularx}\par

23.


\begin{itemize}[leftmargin=2em,itemsep=.2em,topsep=.4em]\item[A.] They are not as willing to help strangers.
\item[B.] They are not as patient with one another.
\item[C.] They violate traffic rules more frequently.
\item[D.] They become more easily irritated in public.
\end{itemize}

24.


\begin{itemize}[leftmargin=2em,itemsep=.2em,topsep=.4em]\item[A.] It was practiced by Boston’s founding fathers.
\item[B.] It is not deemed exotic by Proper Bostonians.
\item[C.] It was adopted by Boston’s upper class.
\item[D.] It is not part of Boston’s local culture.
\end{itemize}

25.


\begin{itemize}[leftmargin=2em,itemsep=.2em,topsep=.4em]\item[A.] Stick to its own way of showing courtesy to strangers.
\item[B.] Follow the examples set by Paris and New York City.
\item[C.] Learn from the world’s major cities in promoting tourism.
\item[D.] Take pride in its history and adhere to its cultural tradition.
\end{itemize}

\subsection*{\textbf{Part III Reading Comprehension (40 minutes)}}


\subsection*{\textbf{Section A}}


\textbf{Directions:} \textit{In this section, there is a passage with ten blanks. You are required to select one word for each blank from a} \textit{list of choices given in a word bank following the passage. Read the passage through carefully before making your} \textit{choices. Each choice in the bank is identified by a letter. Please mark the corresponding letter for each item on }\textbf{\textit{Answer}} \textbf{\textit{Sheet 2}}\textit{ with a single line through the centre. You may not use any of the words in the bank more than once.}


In her second year, Charlene Duong learned of the use of poisonous, synthetic pesticides on her college campus. Shocked but not surprised, she knew she had to do something. Along with a couple of classmates, Charlene did a quick web search and discovered a small but growing movement led by the organization, Herbicide-Free Campus (HFC), to rid college campuses of artificial herbicide (灭草剂). They were \uline{26 }.


Like many, Charlene experiences climate anxiety — a \uline{27} fear of a climate catastrophe — and was, at the time, looking for an \uline{28 }. When she discovered the HFC movement, she said she felt she “had found a specific area to focus on that still fit into the larger picture of fighting for a healthier, safer, cleaner \uline{29} for all.”


Toxic herbicide use in university land care is not unique. Most institutions of higher education rely on synthetic pesticides and fertilizers to achieve \uline{30} goals. Having a “beautiful” campus means green and perfectly maintained lawns along with flower beds and paved sidewalks. But these \uline{31} managed campuses can come at a cost: increased cancer risk, \uline{32} waterways, poisoned wildlife and lifeless soil.


Pesticide use on college campuses also contributes to our global climate crisis. The use of chemicals to get rid of insects or unwanted plant life can increase indirect \uline{33 }, as they can include petroleum-based ingredients. Pesticide use also decreases the life in soil, \uline{34} the ability of soils to absorb carbon or retain water and thus reducing campuses’ ability to recover quickly from climate-related extreme weather events like droughts and floods.


Instead of using toxic chemicals, students working with HFC help out with \uline{35} the campus grounds.


“This work reminds me to be in the present moment as I play my role in reducing herbicide use and keeping my campus safe and healthy,” says Charlene.


\begin{itemize}[leftmargin=2em,itemsep=.2em,topsep=.4em]\item[A.] aesthetic
\item[B.] chronic
\item[C.] contaminated
\item[D.] conventionally
\item[E.] emissions
\item[F.] environment
\item[G.] hampering
\item[H.] incidentally
\item[I.] infringement
\item[J.] intrigued
\item[K.] juvenile
\item[L.] outlet
\item[M.] rotating
\item[N.] vibrations
\item[O.] weeding
\end{itemize}

\subsection*{\textbf{Section B}}


\textbf{Directions:} \textit{In this section, you are going to read a passage with ten statements attached to it. Each statement contains} \textit{information given in one of the paragraphs. Identify the paragraph from which the information is derived. You may} \textit{choose a paragraph more than once. Each paragraph is marked with a letter. Answer the questions by marking the} \textit{corresponding letter on }\textbf{\textit{Answer Sheet 2}}\textit{.}


\textbf{Why Your Library Is the Most Important Place in Town}


A) Librarians know the value of their community services, and their patrons appreciate their importance as well. But in an increasingly digital world, we see the role of libraries as community and cultural centers at times undervalued. When shrinking municipal budgets combine with the nonstop technological revolution, public library services that focus on building communities face-to-face, inspiring and educating patrons about art, literature, and music, and helping patrons engage in civil discourse can seem old-fashioned. But it is precisely those shrinking budgets and the assault of technologically mediated life that make public libraries’ cultural and community offerings more important than ever.


B) Many people point out the value public libraries bring to their communities. More than just books and banks of computers, libraries are still places where individuals gather to explore, interact, and imagine. Some of the specific ways in which libraries add value to our communities and serve as cultural centers for our patrons are community builders, centers for the arts, and champions of youth. Libraries serve in these capacities and are more than just about information.


C) As community builders libraries are engaged in incredible work. From tiny public libraries to huge city institutions with multiple branches, libraries across the United States are building community and supporting local culture in exciting ways. These are inspiring and hopefully will encourage librarians interested in community services and cultural outreach to make connections with each other, share ideas, and build partnerships. Supporters of libraries as community builders claim that unless you are out there changing neighborhoods, you are not completing the work you are to do. Strengthening neighborhoods and championing the cultural lives of communities are big responsibilities.


D) Place-based economic development stresses the importance of offering attractive, functional, and community-based places, such as libraries, in town squares and depressed neighborhoods. Like a major department store in a mall, libraries attract large numbers of people, creating economic opportunities for numerous businesses and organizations in the surrounding area. Large cities, medium-sized ones, and even small towns have successfully transformed their libraries into the hubs (枢纽) of vibrant neighborhoods.


E) As key municipal agencies, and focal points for community education, libraries are major players in creating livable, environmentally friendly cities and towns. The Urban Libraries Council released a report detailing the unique ways in which libraries can further sustainability at the local level. Beyond ensuring that library construction projects consider environmental impact, libraries can take a lead in supporting local foods and craftsmen, like the Peabody Institute Library’s partnering with local businesses to pioneer a farmers’ market in their courtyard, or the Richmond Public Library’s seed lending library which “nurtures locally-adapted plant varieties, and fosters community resilience (韧性), self-reliance and a culture of sharing.”


F) Archives preserve historic artifacts, oral histories, digital history projects, and scholarly writings relevant to the community, including minority groups. Communities lucky enough to have archivists have a great advantage when it comes to organizing historical records and artifacts. An organized archive is a place where people can research their ancestry and immigration history, do environmental research, and more. An archivist is an advocate for preservation who, among other things, coordinates the restoration of maps and paintings, the digitization of vital records, and the creation of oral history projects. With projects like the Massachusetts Memories Road Show and the Veterans History Project, evidence of the importance of archives is everywhere.


G) In the words of Robert Putnam, “People may go to the library looking mainly for information, but they find each other there.” New moms connect at baby story-times; elderly people, often facing difficult life transitions, attend events and find that they make new friends; teenagers meet up in libraries’ teen spaces after school; and readers discuss current events in the periodicals room. In libraries, community-building connections are happening all the time.


H) As Keith Richards said, “The public library is the great equalizer.” Despite the rising costs of concert and theater tickets, public library events (including concerts, author visits, and gallery displays) are often offered free of charge, enabling people of any income level to attend. In addition, library book groups allow people to explore and discuss the literary arts, and the Great Stories Club introduces at-risk youth to literature. The best part: it’s all free and open to the public.


I) In a time when education is increasingly expensive, public libraries provide information and educational opportunities free for all people, regardless of their socio-economic status. Offered by libraries across the country, American Library Association’s Let’s Talk about It programs are wonderful examples of scholar-facilitated learning opportunities in libraries. In addition, many libraries present classes and discussion programs, and some even provide online continuing education courses such as the Universal Class database.


J) Librarians know that patrons aren’t just information consumers, they’re information producers. Patrons use the library to gain knowledge in order to create their own new and independent works. Increasing numbers of libraries provide spaces and services that meet the needs of people who want to learn how to edit Wikipedia, set up blogs or podcasts, create their own magazines, and so much more. Many libraries offer art or writing workshops and groups, and some provide music practice rooms for patrons. Programs like ImaginOn in Charlotte, North Carolina, provide exciting models that take community partnership, creativity, and creation to a new level.


K) The decline of civil discourse stems in part from the fact that it is so easy for people to watch news about, buy products from, and engage — in both the virtual and real worlds — only with those of similar backgrounds and ideologies. Public libraries, through such programs as The Human Library and Socrates Café, can help build small communities of difference that encourage people to interact with and learn from each other through dialogue. By both actively promoting civil discourse through these programs, and modeling and upholding the principles of free inquiry and expression for all, libraries help individuals rediscover the importance of and increased need for civil discourse in American life.


L) Free tutoring, homework help programs, and summer reading programs for kids and teens help bridge the economic divide that impacts students’ academic performance. The cost of hiring a private tutor is well beyond what many library patrons can afford, so libraries offer homework help and tutoring online, by phone, in person, and even through social media and homework apps. Annual summer reading programs also have a positive impact on student performance and, according to a 2010 study conducted by Dominican University’s Graduate School of Library and Information Science, students’ reading skills get a boost from these popular nationwide events.


M) Through library collections, programs, and physical spaces, children learn to share, to be engaged in their communities, to participate in the arts, and to explore their immediate world and the world at large. There are surely endless examples of innovative library services for children, including the Middle Country Public Library’s Nature Explorium, which engages children in learning about the natural world.


N) These examples are just a few of the many amazing things that public libraries around the United States (and the world) are doing to build and maintain strong community connections. We encourage you to try some of these ideas in your own libraries, and we hope that these ideas will help you be better able to convince your community leaders of the important role that public libraries play in communities large and small.


36. People going to the library in search of information can build connections with each other there.


37. According to advocates of libraries as community builders, librarians are not doing their job well if they do not


change their communities.


38. With the costs of education continually rising, public libraries remain places where all people can have access to


education.


39. Libraries draw large crowds, thus creating lots of business opportunities in neighboring areas.


40. With the world more and more digitalized, people sometimes underestimate the role of libraries as community and


cultural centers.


41. Various programs organized by public libraries for children and adolescents help narrow the gap between students


from varying economic backgrounds.


42. In an organized archive, people can do research on their family history and find out how their ancestors came to


settle in the new land.


43. Public libraries organize cultural events, often allowing people of different income levels to attend free of charge.


44. Besides being an information provider, the library performs many other important services for the community.


45. Public libraries can help build small communities of people with different backgrounds and ideologies.


\subsection*{\textbf{Section C}}


\textbf{Directions: }\textit{There are 2 passages in this section. Each passage is followed by some questions or unfinished statements.} \textit{For each of them there are four choices marked A, B, C and D. You should decide on the best choice and mark the} \textit{corresponding letter on }\textbf{\textit{Answer Sheet 2}}\textit{ with a single line through the centre.}


\subsection*{\textbf{Passage One}}


\textbf{Questions 46 to 50 are based on the following passage.}


Nationally, one in six children miss 15 or more days of school in a year. Education officials have deplored all this missed instruction.


These chronically absent students suffer academically because of all the classroom instruction they miss out on. In 2015, the US secretary of education responded to this crisis, urging communities to support every student to attend every day and be successful in school. His open letter stated that missing 10\% of school days in a year for any reason — excused or unexcused — is a primary cause of low academic achievement.


Worrying about whether children attend school makes sense. After all, if students don’t show up, teachers can’t teach them.


But what if America’s attendance crisis is about much more than students missing class? What if, instead, it is a reflection of family and community crises these students face — such as being ejected from the family apartment, fearing for their safety in their neighborhood or suffering an illness?


As social scientists we investigated how excused and unexcused absences relate to children’s academic achievement.


We find that absences excused by a parent do little to harm children’s learning. In fact, children with no unexcused absences — but 15 to 18 excused absences — have test scores equal to their peers who have no absences.


Meanwhile, the average child with even just one unexcused absence does much worse academically than peers with none.


We believe unexcused absence is a strong signal of the many challenges children and families face, including economic and medical hardships. Unexcused absences can be a powerful signal of how those out-of-school challenges affect children’s academic progress.


Our evidence suggests unexcused absences are problematic, but for a different reason than people often think. Absence from school, and especially unexcused absence, matters mainly as a signal of many crises children and their families may be facing. It matters less as a cause of lower student achievement due to missed instruction.


How we choose to think of school absences matters for educational policy. School attendance policies typically hold schools and families accountable for the days children miss, regardless of whether they were excused or unexcused absences.


These policies assume that missing school for any reason harms children academically because they are missing classroom instruction. They also assume that schools will be able to effectively intervene by reducing student absences. We find neither to be the case.


As a result, these attendance policies end up disproportionately punishing families dealing with out-of-school crises in their lives and pressuring schools who serve them to get students to school more often.


We instead suggest using unexcused absence from school as a signal to channel resources to the children and families who need them most.


46. What does the US secretary of education say in his open letter?


\begin{itemize}[leftmargin=2em,itemsep=.2em,topsep=.4em]\item[A.] It is of vital importance to respond promptly to the school absence crisis.
\item[B.] The academic performance of chronically absent students is deplorable.
\item[C.] Low academic achievement is mainly attributed to school absences.
\item[D.] The effect of school absences on American education is worrisome.
\end{itemize}

47. What do the authors find about school absences?


\begin{itemize}[leftmargin=2em,itemsep=.2em,topsep=.4em]\item[A.] Excused school absences have little impact on children’s learning.
\item[B.] There is little difference between unexcused and excused absences.
\item[C.] Excused absences lead to comparatively better school performance.
\item[D.] Unexcused absences are a big challenge to both schools and families.
\end{itemize}

48. What do the authors believe concerning unexcused school absences?


\begin{itemize}[leftmargin=2em,itemsep=.2em,topsep=.4em]\item[A.] They are likely to cause a decrease in students’ academic achievements due to missed instruction.
\item[B.] They point directly to many of the out-of-school challenges confronting children and their families.
\item[C.] They are matters the American government typically ignores when formulating educational policies.
\item[D.] They give a clear signal to children and their families of the crises they are likely to face in the future.
\end{itemize}

49. What is the assumption underlying education policies in the US?


\begin{itemize}[leftmargin=2em,itemsep=.2em,topsep=.4em]\item[A.] Children’s academic performance depends on reducing the number of absences.
\item[B.] Schools can boost children’s academic performance by effective intervention.
\item[C.] Schools as well as families should be held responsible for out-of-school crises.
\item[D.] Children’s academic performance is closely related to the quality of instruction.
\end{itemize}

50. What do the authors suggest doing regarding school absences?


\begin{itemize}[leftmargin=2em,itemsep=.2em,topsep=.4em]\item[A.] Identifying their underlying causes.
\item[B.] Reframing school attendance policies.
\item[C.] Directing resources to helping needy children.
\item[D.] Pressuring schools to reduce unexcused ones.
\end{itemize}

\subsection*{\textbf{Passage Two}}


\textbf{Questions 51 to 55 are based on the following passage.}


After earning a bachelor’s degree, I was determined to do what I love. I headed straight to graduate school to investigate the social problems that fascinated me. For almost a decade, I told everyone I encountered that they should do the same. “Follow your passion,” I counseled. “You can figure out the employment stuff later.” It wasn’t until I began to research this widely accepted career advice that I understood how problematic it really was. As a sociologist, I interviewed college students and professional workers to learn what it really meant to pursue their dreams, which I will refer to here as the passion principle. I was stunned by what I found out about this principle in the research for my new book. Surveys show the American public has long held the passion principle in high regard as a career decision-making priority. And its popularity is even stronger among those facing job instability. Advocates of the passion principle found it compelling because they believed that following one’s passion can provide workers with both the motivation necessary to work hard and a place to find fulfillment. Yet, what I found is that following one’s passion does not necessarily lead to fulfillment, but is one of the most powerful cultural forces perpetuating overwork. I also found that promoting the pursuit of one’s passion helps perpetuate social inequalities due to the fact that not everyone has the same economic resources to allow them to pursue their passion with ease.


While the passion principle is broadly popular, not everyone has the necessary resources to turn their passion into a stable, good-paying job. Passion-seekers from wealthy families are better able to wait until a job they are passionate about comes along without worrying about student loans in the meantime. They are also better situated to take unpaid internships to get their foot in the door while their parents pay their rent. And they often have access to parents’ social networks to help them find jobs. Surveys revealed that working-class and first-generation college graduates, regardless of their career field, are more likely than their wealthier peers to end up in low-paying unskilled jobs when they pursue their passion.


Colleges, workplaces and career counselors who promote the “follow your passion” path for everyone, without leveling the playing field, help perpetuate socioeconomic inequalities among career aspirants.


It’s not just well-off passion-seekers who benefit from the passion principle. Employers of passionate workers do, too. Potential employers showed greater interest in passionate applicants in part because they believed the applicants would work hard at their jobs without expecting an increase in pay. They even sacrifice a good salary, job stability and leisure time to work in a job they love.


51. What did the author advise people to do for almost a decade?


\begin{itemize}[leftmargin=2em,itemsep=.2em,topsep=.4em]\item[A.] Figure out what is the most fascinating job.
\item[B.] Follow widely accepted career counsel.
\item[C.] Pursue their careers with passion.
\item[D.] Do whatever they are zealous for.
\end{itemize}

52. How did the author feel about the passion principle through his research?


\begin{itemize}[leftmargin=2em,itemsep=.2em,topsep=.4em]\item[A.] He was astonished by its consequences.
\item[B.] He was further convinced of its soundness.
\item[C.] He was actually right to follow it through.
\item[D.] He was struck by its broad popularity.
\end{itemize}

53. What is important to turning one’s passion into a stable, good-paying job?


\begin{itemize}[leftmargin=2em,itemsep=.2em,topsep=.4em]\item[A.] Willingness to take unpaid internships and low-paying jobs.
\item[B.] Full academic preparedness and sound career counseling.
\item[C.] Hard work and sacrifice of leisure time.
\item[D.] Financial backing and social connections.
\end{itemize}

54. What happens when everyone is encouraged to follow their passion?


\begin{itemize}[leftmargin=2em,itemsep=.2em,topsep=.4em]\item[A.] Many more career aspirants end up unemployed.
\item[B.] People are less concerned with socioeconomic inequality.
\item[C.] Socioeconomic inequality is likely to persist.
\item[D.] Career counselors are going to lose credibility.
\end{itemize}

55. What does the author say about employers of passionate workers?


\begin{itemize}[leftmargin=2em,itemsep=.2em,topsep=.4em]\item[A.] They provide these workers with job stability and a good salary.
\item[B.] They exploit these workers’ passion to benefit themselves.
\item[C.] They level the playing field for these workers to reach their goals.
\item[D.] They encourage these workers to realize their aspirations.
\end{itemize}

\subsection*{\textbf{Part IV Translation} \textbf{(30 minutes)}}


\textbf{Directions: }\textit{For this part, you are allowed 30 minutes to translate a passage from Chinese into English. You should} \textit{write your answer on }\textbf{\textit{Answer Sheet 2}}\textit{.}


天宫空间站 (Tiangong Space Station) 是中国首个太空实验室，拥有110 多立方米使用空间，可驻留3 名航天员，在距地球表面400—450 公里的轨道上运行。天宫空间站已实施180 多个科学研究与应用项目，涉及空间生命科学、太空医学、空间材料科学等领域。天宫空间站的研究成果在我国得到了广泛应用，产生了显著的经济效益。例如，太空育种创造的直接经济效益高达数千亿元。这不仅标志中国在航天技术上取得了巨大进步， 也表明中国为全球的太空研究和应用做出了重大贡献。

\end{document}